% =====================================================================
% 设计说明书 LaTeX 模板
% 使用说明：将 {{变量名}} 替换为实际内容
% 截图：使用 \screenshot{文件名}{标题} 宏
% 编译：xelatex -interaction=nonstopmode 文件名.tex （运行两遍）
% =====================================================================
\documentclass[a4paper,12pt]{article}
\usepackage[UTF8,heading=true]{ctex}
\usepackage[top=2.5cm,bottom=2.5cm,left=2.5cm,right=2.5cm]{geometry}
\usepackage{fancyhdr}
\usepackage{lastpage}
\usepackage{graphicx}
\usepackage{tabularx}
\usepackage{enumitem}
\usepackage{hyperref}
\usepackage{xcolor}
\usepackage{float}
\usepackage{longtable}
\usepackage{array}

\hypersetup{hidelinks}
\graphicspath{{images/}}

\ctexset{
  section/format = \Large\bfseries,
  subsection/format = \large\bfseries,
  subsubsection/format = \normalsize\bfseries,
}

\pagestyle{fancy}
\fancyhf{}
\fancyhead[C]{{{软件全称}}V{{版本号}}\quad 设计说明书}
\fancyfoot[C]{第 \thepage\ 页\quad 共 \pageref{LastPage} 页}
\renewcommand{\headrulewidth}{0.5pt}
\renewcommand{\footrulewidth}{0.5pt}
\fancypagestyle{plain}{\pagestyle{fancy}}

\setlength{\parindent}{2em}
\setlength{\parskip}{0.3em}
\renewcommand{\baselinestretch}{1.3}

\newcommand{\screenshot}[2]{%
  \begin{figure}[H]
    \centering
    \includegraphics[width=0.35\textwidth]{#1}
    \caption{#2}
  \end{figure}
}

\begin{document}

\begin{center}
  {\LARGE\bfseries {{软件全称}}\quad 设计说明书}

  \vspace{0.8em}
  {\large 著作权人：{{著作权人名称}}}

  \vspace{0.3em}
  软件版本：V{{版本号}}\qquad 发布日期：{{首次发表日期}}
\end{center}

\vspace{1.5em}

% =================================================================
\section{软件概述}
% =================================================================

\subsection{软件简介}

{{软件简介}}

\subsection{开发背景}

{{开发背景，描述行业痛点和用户需求}}

\subsection{设计目标}

\begin{itemize}[leftmargin=2em]
  \item {{目标1}}
  \item {{目标2}}
  \item {{目标3}}
\end{itemize}


% =================================================================
\section{需求分析}
% =================================================================

\subsection{功能需求}

\begin{enumerate}[leftmargin=2em]
  \item {{功能需求1}}。
  \item {{功能需求2}}。
  \item {{功能需求3}}。
\end{enumerate}

\subsection{性能需求}

\begin{itemize}[leftmargin=2em]
  \item 响应时间：{{响应时间要求}}。
  \item 并发支持：{{并发要求}}。
  \item 资源占用：{{资源要求}}。
\end{itemize}

\subsection{安全需求}

\begin{itemize}[leftmargin=2em]
  \item 数据加密：{{加密方案}}。
  \item 权限控制：{{权限方案}}。
  \item 防护措施：{{安全措施}}。
\end{itemize}


% =================================================================
\section{总体设计}
% =================================================================

\subsection{系统架构}

{{系统架构描述，前后端架构、数据流}}

\subsection{模块划分}

软件主要包含以下模块：

\begin{itemize}[leftmargin=2em]
  \item {{模块1名称}}（{{对应目录}}）：{{模块1功能简述}}。
  \item {{模块2名称}}（{{对应目录}}）：{{模块2功能简述}}。
  \item {{模块3名称}}（{{对应目录}}）：{{模块3功能简述}}。
\end{itemize}

\subsection{技术选型}

\begin{itemize}[leftmargin=2em]
  \item {{技术1}}。
  \item {{技术2}}。
  \item {{技术3}}。
\end{itemize}

\subsection{运行环境}

\begin{itemize}[leftmargin=2em]
  \item \textbf{平台}：{{运行平台}}。
  \item \textbf{客户端}：{{客户端要求，列出所有支持的操作系统}}。
  \item \textbf{基础库版本}：{{版本信息}}。
\end{itemize}


% =================================================================
\section{详细设计}
% =================================================================

% ⚠️ 强制规则（避免补正）：
%   1. 每个在「源代码鉴别材料（前30页+后30页）」中出现过的模块/目录，
%      在本节都必须有对应的 \subsection；反之，本节出现的每个模块，
%      都必须能在源程序文档里看到对应代码。审查员会做双向核对。
%   2. 每个 \subsection 标题必须与项目实际目录或文件名对应（如
%      pages/index、utils/api.js），不要起架空的中文名。
%   3. 每个模块按下面的 6 段结构展开，缺段会被打回。

% --- 为每个核心模块重复以下 6 段结构 ---

\subsection{{{模块名称（与源代码目录/文件对应）}}}

\subsubsection{模块职责}

{{该模块在系统中承担的核心职责，一句话定位 + 1--2 句展开说明}}

\subsubsection{输入与输出}

\begin{itemize}[leftmargin=2em]
  \item \textbf{输入}：{{用户操作 / 上游模块传入的数据 / 跳转参数 / 接口入参}}。
  \item \textbf{输出}：{{渲染结果 / 下游模块消费的数据 / 跳转目标 / 接口返回}}。
\end{itemize}

\subsubsection{关键算法与处理逻辑}

{{核心算法或业务逻辑描述，例如距离计算、日期匹配、过滤排序、状态转换等。
若涉及公式或判定规则，用列表或简短伪代码展示}}

\subsubsection{状态与缓存管理}

\begin{itemize}[leftmargin=2em]
  \item \textbf{本地状态}：{{页面/组件 data 字段、内存变量}}。
  \item \textbf{持久化}：{{localStorage / 数据库 / 全局 globalData 等}}。
  \item \textbf{失效与刷新策略}：{{何时清空、何时同步}}。
\end{itemize}

\subsubsection{与其他模块的依赖关系}

\begin{itemize}[leftmargin=2em]
  \item \textbf{依赖}：{{本模块调用了哪些上游模块/工具/接口}}。
  \item \textbf{被依赖}：{{哪些下游模块依赖本模块的输出}}。
\end{itemize}

\subsubsection{处理流程}

\begin{enumerate}[leftmargin=2em]
  \item {{流程步骤1}}。
  \item {{流程步骤2}}。
  \item {{流程步骤3}}。
\end{enumerate}

% \screenshot{图-模块截图.png}{模块界面}

% --- 重复以上模式 ---


% =================================================================
\section{数据结构设计}
% =================================================================

\subsection{数据存储}

\begin{itemize}[leftmargin=2em]
  \item \textbf{本地存储}：{{本地存储方案和内容}}。
  \item \textbf{云端数据库}：{{云端存储方案和内容}}。
\end{itemize}

\subsection{数据格式}

\begin{itemize}[leftmargin=2em]
  \item {{数据格式1描述}}。
  \item {{数据格式2描述}}。
\end{itemize}

\subsection{数据关系}

\begin{itemize}[leftmargin=2em]
  \item {{数据关系1}}。
  \item {{数据关系2}}。
\end{itemize}


% =================================================================
\section{接口设计}
% =================================================================

\subsection{内部接口}

\begin{itemize}[leftmargin=2em]
  \item {{内部接口1描述}}。
  \item {{内部接口2描述}}。
\end{itemize}

\subsection{外部接口}

\begin{center}
\begin{tabularx}{\textwidth}{|l|l|X|}
\hline
\textbf{接口} & \textbf{方法} & \textbf{说明} \\
\hline
{{接口路径1}} & {{方法1}} & {{说明1}} \\
\hline
{{接口路径2}} & {{方法2}} & {{说明2}} \\
\hline
\end{tabularx}
\end{center}


% =================================================================
\section{界面设计}
% =================================================================

\subsection{整体布局}

{{整体布局描述}}

% \screenshot{图-整体布局.png}{整体布局}

\subsection{主要页面}

% --- 为每个页面添加描述和截图 ---

\subsubsection{{{页面名称}}}

{{页面描述}}

% \screenshot{图-页面.png}{页面说明}

% --- 重复以上模式 ---

\subsection{交互设计}

% --- 描述主要交互方式 ---


% =================================================================
\section{安全设计}
% =================================================================

\subsection{认证机制}

\begin{itemize}[leftmargin=2em]
  \item {{认证方式1}}。
  \item {{认证方式2}}。
\end{itemize}

\subsection{数据保护}

\begin{itemize}[leftmargin=2em]
  \item {{数据保护措施1}}。
  \item {{数据保护措施2}}。
\end{itemize}

\subsection{访问控制}

\begin{itemize}[leftmargin=2em]
  \item {{访问控制1}}。
  \item {{访问控制2}}。
\end{itemize}


% =================================================================
\section{测试设计}
% =================================================================

\subsection{测试策略}

\begin{itemize}[leftmargin=2em]
  \item 功能测试：{{功能测试范围}}。
  \item 兼容性测试：{{兼容性测试范围}}。
  \item 性能测试：{{性能测试范围}}。
\end{itemize}

\subsection{测试环境}

\begin{itemize}[leftmargin=2em]
  \item {{测试环境1}}。
  \item {{测试环境2}}。
\end{itemize}

\end{document}
